\documentclass[11pt,a4paper]{article}
\usepackage[utf8]{inputenc}
\usepackage[T1]{fontenc}
\usepackage[spanish]{babel}
\usepackage{lmodern}
\usepackage[margin=2.5cm]{geometry}
\usepackage[hidelinks]{hyperref}
\setlength{\parindent}{0pt}
\setlength{\parskip}{0.6em}
\pagestyle{plain}
\begin{document}
{\small
\begin{tabular}{@{}l@{}}
\textbf{Curso: Matemática Discreta} \\
\textbf{Sección: 30} \\
Ivan Morataya --- 16667 \\
David Díaz --- 25314 \\
David Ibaté --- 251732 \\
Miguel Carranza --- 24458 \\
José Ceferino --- 251043
\end{tabular}\par
}
\vspace{0.5em}
{\large\bfseries Proyecto 3: Permutaciones y combinaciones}\par
Guía para entender y utilizar el programa


\section*{1. Descripción del problema}
Este programa responde dos preguntas: ¿de cuántas maneras podemos ordenar
unos objetos? y ¿de cuántas maneras podemos elegir algunos de ellos?
Además de indicar la cantidad, muestra las posibilidades en una lista.

Los objetos pueden ser letras, números o nombres. Por ejemplo, las letras
\texttt{a b c d} representan cuatro objetos diferentes. La lista
\texttt{a a b c} representa cuatro objetos, pero dos son iguales.

El programa se utiliza escribiendo opciones en un menú. No es necesario
modificar el código ni saber programar.

\subsection*{¿Qué es una permutación?}
Es una manera de \textbf{ordenar todos los objetos}. El orden sí importa.
Con las letras \texttt{a b}, los órdenes \texttt{(a, b)} y
\texttt{(b, a)} son dos resultados diferentes.

\subsection*{¿Qué es una combinación?}
Es una manera de \textbf{elegir una cantidad de objetos}. El orden no importa.
Elegir \texttt{a} y \texttt{b} es lo mismo que elegir \texttt{b} y
\texttt{a}; por eso se cuenta una sola vez.

\subsection*{Las cuatro opciones del programa}
\begin{enumerate}
\item \textbf{Permutaciones de objetos diferentes:} ordenar todos los
objetos cuando ninguno se repite. Ejemplo: \texttt{a b c d}.
\item \textbf{Permutaciones de objetos iguales:} ordenar todos los objetos
cuando algunos se repiten. Ejemplo: \texttt{a a b c}. Intercambiar las dos
letras \texttt{a} entre sí no crea un orden nuevo.
\item \textbf{Combinaciones de objetos diferentes:} elegir algunos objetos
de una lista sin repetidos. Ejemplo: elegir 2 de \texttt{a b c d}.
\item \textbf{Combinaciones de objetos iguales:} elegir algunos objetos
de una lista que puede tener repetidos. Ejemplo: elegir 2 de
\texttt{a a b c}. Solo se pueden usar las copias disponibles.
\end{enumerate}

\newpage
\section*{2. Cómo utilizar el programa}
\subsection*{Abrir el menú}
Se necesita tener instalado Python 3.10 o posterior y contar con el archivo
\texttt{proyecto.py}. El programa no requiere instalar otros paquetes.

\begin{enumerate}
\item Abre una terminal: es la ventana en la que se escriben comandos.
En este equipo puedes abrirla con \textbf{Ctrl + Alt + T}.
\item En este equipo, escribe lo siguiente y presiona \textbf{Enter}:
\begin{verbatim}
cd ~/Documentos/MAD/VIS/Proyecto-3-Discreto
\end{verbatim}
Este comando abre la carpeta del proyecto. Si guardas el proyecto en otra
carpeta, utiliza la ruta de esa carpeta.
\item Escribe este comando y presiona \textbf{Enter}:
\begin{verbatim}
python3 proyecto.py --interactivo
\end{verbatim}
\item Aparecerá un menú con las opciones 1, 2, 3 y 4 descritas en la página
anterior. La opción \textbf{5} ejecuta los cuatro ejemplos del proyecto y
la opción \textbf{0} cierra el programa.
\end{enumerate}

\subsection*{Hacer una consulta, paso a paso}
Para elegir dos objetos de la lista \texttt{a a b c}:
\begin{enumerate}
\item En \texttt{Opción:}, escribe \texttt{4} y presiona Enter.
\item Cuando solicite los objetos, escribe \texttt{a a b c} y presiona Enter.
Separa los objetos con espacios, sin comas.
\item En \texttt{Cantidad a seleccionar r:}, escribe \texttt{2} y presiona Enter.
La letra \texttt{r} significa cuántos objetos quieres elegir.
\item El programa mostrará \texttt{Total: 4} y esta lista:
\begin{verbatim}
1. (a, a)
2. (a, b)
3. (a, c)
4. (b, c)
\end{verbatim}
\item El menú vuelve a aparecer. Puedes hacer otra consulta o escribir
\texttt{0} para salir.
\end{enumerate}

\subsection*{Cómo escribir los datos y leer la respuesta}
Cada palabra separada por un espacio cuenta como un objeto. Por ejemplo,
\texttt{rojo azul verde} son tres objetos. Escribe \texttt{a a} para indicar
dos copias iguales; \texttt{aa} sin espacio representa un solo objeto.
El programa distingue \texttt{a} de \texttt{A}.

\texttt{Total} indica cuántas posibilidades existen. Los paréntesis agrupan
los objetos de cada resultado. En las opciones 1 y 2 se usan todos los
objetos; en las opciones 3 y 4 se pregunta cuántos quieres elegir.

\newpage
\section*{3. Casos de prueba y resultados obtenidos}
Cada caso se ejecutó en el programa y se comparó con el cálculo esperado.
En las listas de esta página, \texttt{abcd} significa el orden
\texttt{(a, b, c, d)}. Se escribe abreviado para facilitar la lectura.

\subsection*{Caso C1. Ordenar objetos diferentes}
\textbf{Entrada:} \texttt{a b c d}. \textbf{Opción del menú:} 1.

Hay 4 opciones para el primer lugar, 3 para el segundo, 2 para el tercero y
1 para el último. Por tanto, $4\times3\times2\times1=24$.

\textbf{Esperado: 24. Obtenido: 24. Estado: correcto.}
\begin{verbatim}
abcd  abdc  acbd  acdb  adbc  adcb
bacd  badc  bcad  bcda  bdac  bdca
cabd  cadb  cbad  cbda  cdab  cdba
dabc  dacb  dbac  dbca  dcab  dcba
\end{verbatim}

\subsection*{Caso C2. Ordenar objetos con repetidos}
\textbf{Entrada:} \texttt{a a b c}. \textbf{Opción del menú:} 2.

Las dos letras \texttt{a} son iguales. Si se contaran como diferentes,
cada orden aparecería dos veces. Se corrige dividiendo entre 2:
$24\div2=12$.

\textbf{Esperado: 12. Obtenido: 12. Estado: correcto.}
\begin{verbatim}
aabc  aacb  abac  abca  acab  acba
baac  baca  bcaa  caab  caba  cbaa
\end{verbatim}

\subsection*{Caso C3. Elegir objetos diferentes}
\textbf{Entrada:} \texttt{a b c d}. \textbf{Opción:} 3.
\textbf{Cantidad a elegir:} 2.

Hay $4\times3=12$ parejas si se cuenta el orden. Como cada pareja aparece
en dos órdenes, se divide entre 2: $12\div2=6$.

\textbf{Esperado: 6. Obtenido: 6. Estado: correcto.}
\begin{verbatim}
(a, b)  (a, c)  (a, d)  (b, c)  (b, d)  (c, d)
\end{verbatim}

\subsection*{Caso C4. Elegir objetos con repetidos}
\textbf{Entrada:} \texttt{a a b c}. \textbf{Opción:} 4.
\textbf{Cantidad a elegir:} 2.

\textbf{Esperado: 4. Obtenido: 4. Estado: correcto.}
\begin{verbatim}
(a, a)  (a, b)  (a, c)  (b, c)
\end{verbatim}
Se permite \texttt{(a, a)} porque hay dos copias de \texttt{a}. No se
permiten \texttt{(b, b)} ni \texttt{(c, c)} porque solo hay una copia de
cada una. Cambiar el orden de una pareja no crea otra combinación.

\newpage
\section*{4. Pruebas adicionales}
También se revisó qué sucede con objetos iguales, selecciones vacías y datos
incorrectos. En todos los casos de la tabla, el resultado obtenido coincidió
con el esperado. Las opciones son las del menú de la página 2.

{\small
\renewcommand{\arraystretch}{1.35}
\begin{tabular}{p{0.9cm}p{4.0cm}p{3.8cm}p{3.8cm}}
\hline
\textbf{Caso} & \textbf{Entrada} & \textbf{Esperado} & \textbf{Obtenido} \\
\hline
C5 & Opción 2: a a a & 1 orden: (a,a,a) & 1 orden: (a,a,a) \\
C6 & Opción 4: a a b; elegir 2 & 2 selecciones: (a,a) y (a,b) & Las mismas 2 selecciones \\
C7 & Opción 3: a b; elegir 0 & 1 selección vacía: () & 1 selección: () \\
C8 & Opción 4: lista vacía; elegir 0 & 1 selección vacía: () & 1 selección: () \\
C9 & Opción 3: a a b; elegir 2 & Error por objetos repetidos & Error por objetos repetidos \\
C10 & Opción 4: a a b; elegir 4 & Error: solo hay 3 objetos & Error: r debe estar entre 0 y 3 \\
C11 & Opción 4: a a b; elegir $-1$ & Error por cantidad negativa & Error: r debe estar entre 0 y 3 \\
C12 & Ordenar a b c d; límite de visualización 2 & Total 24; mostrar solo 2 & Total 24; 2 visibles y aviso \\
C13 & Opción 2: rojo azul rojo & 3 órdenes diferentes & 3 órdenes sin repetir \\
\hline
\end{tabular}
}

\subsection*{Comprobación del menú}
Se ejecutaron las cuatro opciones del menú con los datos de C1 a C4. También
se introdujo una cantidad incorrecta y luego una correcta: el programa mostró
el error y permitió hacer otra consulta sin cerrarse. Al escribir una palabra
en lugar de la cantidad, también informó el error y regresó al menú.

\subsection*{Evidencia de las pruebas}
Se ejecutaron diez pruebas automatizadas, que son comprobaciones repetibles
del funcionamiento. Todas terminaron correctamente. Dentro de ellas se
compararon los resultados con otra forma de enumerarlos para las 364 listas
pequeñas de longitud 0 a 5 formadas con a, b y c.

Los registros completos están en \texttt{pruebas\_automatizadas.txt},
\texttt{casos\_verificados.txt} y \texttt{resultados.txt}, dentro de
\texttt{documentacion/}. Para repetir las pruebas desde la carpeta principal:
\begin{verbatim}
python3 -m unittest discover -s tests -v
\end{verbatim}
El resultado de la ejecución fue:
\begin{verbatim}
Ran 10 tests
OK
\end{verbatim}

\newpage
\section*{5. Recomendaciones y mensajes del programa}
\begin{itemize}
\item En las opciones de objetos diferentes, no repitas etiquetas.
Si necesitas escribir \texttt{a a b c}, utiliza la opción 2 o la 4.
\item La cantidad a elegir debe ser un número entero entre 0 y el total
de objetos disponibles. No es posible elegir 5 objetos de una lista de 4.
Elegir 0 produce una sola selección vacía, escrita como \texttt{()}.
\item Si aparece \texttt{Error}, lee el mensaje y vuelve a ingresar los
datos cuando aparezca el menú. Por ejemplo, elige la opción 4 otra vez
para corregir el tamaño de una combinación.
\item El programa acepta hasta 100 objetos por consulta. Para comenzar,
es conveniente utilizar listas pequeñas.
\item Por defecto se muestran como máximo 50 resultados. Si hay más, el
programa lo indica. El número que aparece en \texttt{Total} sigue siendo
la cantidad completa. Los cuatro ejemplos de este informe sí se muestran
completos.
\end{itemize}

\section*{6. Metodología utilizada}
Primero se distinguió si el problema pide ordenar o elegir. Después se
revisó si hay objetos repetidos. Con esas dos decisiones se identificaron
los cuatro casos que ofrece el menú.

El programa cuenta las posibilidades y genera una lista con cada resultado.
Cuando encuentra objetos iguales, lleva el control de cuántas copias quedan
disponibles. En las combinaciones evita contar dos veces una misma selección
escrita en distinto orden.

Para revisar su funcionamiento se ejecutaron los cuatro ejemplos y se
compararon sus listas con los resultados esperados. También se probaron
datos incorrectos, listas vacías y otras listas pequeñas. Las diez pruebas
automatizadas terminaron correctamente. Los totales obtenidos fueron
\textbf{24, 12, 6 y 4}, respectivamente.

\section*{7. Conclusiones}
El programa permite contar y consultar formas de ordenar o elegir objetos
sin hacer todas las listas a mano. La decisión principal es saber si el
orden importa y si existen objetos repetidos.

Para utilizarlo basta con abrir el menú, elegir una opción e ingresar los
datos. En los casos con objetos iguales se respetan las copias disponibles;
no se agregan objetos que no estén en la lista inicial.

\section*{Referencias}
Rosen, K. H. (2019). \textit{Discrete Mathematics and Its Applications}
(8.\textsuperscript{a} ed.). McGraw-Hill Education. Secciones 6.3 (pp. 428--436), 6.5 (pp. 445--456) y 6.6 (pp. 457--460).

Python Software Foundation. \textit{Documentación de Python 3}: módulos itertools, collections, math y unittest.
\url{https://docs.python.org/es/3/}.
\end{document}
